%% kapa: exam=1 round=20 year=2009 period=1 form=A subject=civil_law subject_ko=민법 seq=17 number=17
%% source: Q-Net 감정평가사 기출문제 (cst00309) artlSeq=5024178 file=extracted/1008/5024178/2009년 제20회 감정평가사 1차시험 기출문제/1교시.hwp
%% conversion: kordoc 4.24.51 -> tools/gen1.py (auto); review=proofread-sonnet; flags=none
\begin{question}{17}
주물ㆍ종물에 관한 설명으로 옳지 \uline{않은} 것은? (다툼이 있으면 판례에 의함)
\begin{choices}
  \item 주물의 소유자나 이용자의 상용에 공여되고 있다면 주물 자체의 효용과 직접적 관련이 되지 않는 물건도 종물이다.
  \item 주물의 구성부분은 종물이 아니다.
  \item 종물은 주물의 처분에 따른다는 규정은 임의규정이다.
  \item 건물에 대한 저당권이 실행되어 매수인(경락인)이 건물의 소유권을 취득한 때에는 특별한 사정이 없는 한 그 건물의 소유를 목적으로 한 토지의 임차권도 매수인에게 이전된다.
  \item 주물의 상용에 이바지하는지 여부와 관련하여 주물과 종물 사이에 밀접한 장소적 관련성이 있어야 한다.
\end{choices}
\end{question}
\answer{1}
